\chapter{Your Toolkit and Environment Setup}
\label{ch:toolkit}

\chapterguide%
  {Basic familiarity with files and folders is enough; no previous R installation is assumed.}
  {install R and RStudio, optionally manage R through Anaconda/conda, configure Git and GitHub with SSH, create a reproducible course project, verify the software stack, and understand the datasets used throughout the course.}

\section{What you are installing---and why}

The course uses \textbf{R} as the statistical computing engine and \textbf{RStudio} as the main integrated development environment (IDE). These are related but separate pieces of software. R performs the computation; RStudio provides an editor, console, data viewer, plotting pane, help system, project management, and an integrated terminal.

\begin{utpfig}[The course software stack]{R computes, RStudio provides the working environment, and the operating system supports both. Git records changes to course code; pushing shares committed changes with the GitHub repository.}{fig:r-stack}
\begin{tikzpicture}[node distance=0.35cm,
  layer/.style={draw=UTPBlue!70,fill=UTPPaleBlue,rounded corners=1.5mm,
    line width=0.7pt,text width=0.90\linewidth,minimum height=1.40cm,
    inner xsep=3.5mm,inner ysep=3mm,align=left,
    font=\normalsize\color{UTPNavy}},
  flow/.style={-{Stealth[length=2mm,width=1.5mm]},
    line width=0.8pt,draw=UTPTeal}]
\node[layer] (project) {\textbf{Course project --- files and evidence}\\[3pt]
 \texttt{all\_datasets/}, and \texttt{all\_experiments/}};
\node[layer,below=of project,draw=UTPTeal!80,fill=UTPPaleTeal] (studio)
  {\textbf{RStudio Desktop --- working environment}\\[3pt]
   Edit scripts; use the console, plots, help, and project tools};
\node[layer,below=of studio] (r)
  {\textbf{R --- computation engine}\\[3pt]
   Work with data, fit statistical models, and create graphics};
\node[layer,below=of r,draw=UTPMuted!70,fill=UTPPaleGray] (os)
  {\textbf{Operating system --- platform}\\[3pt]
   Windows, macOS, or Linux};
\node[layer,below=of os,draw=UTPOrange!80,fill=UTPPaleOrange] (git)
  {\textbf{Git and GitHub --- version and share code}\\[3pt]
   Commit reviewed changes, then push to the remote repository.};
\draw[flow] (project.south) -- (studio.north);
\draw[flow] (studio.south) -- (r.north);
\draw[flow] (r.south) -- (os.north);
\draw[flow,dashed] (os.south) -- (git.north);
\end{tikzpicture}
\end{utpfig}

\begin{keyidea}
Install \textbf{R first}, then install \textbf{RStudio}. RStudio does not replace R; it discovers and launches an installed R version.
\end{keyidea}

\section{Course project layout}
Use \texttt{UTP-DATA-SCIENCE/} as the project root, with these folders:

\noindent
\begin{minipage}[t]{0.49\linewidth}
\begin{lstlisting}[style=utpterminal]
all_datasets/
  ames_housing_dataset.csv
  bike_sharing_dataset_day.csv
  bike_sharing_dataset_hour.csv
  breast_cancer_dataset.csv
  credit_card_fraud_dataset.csv
  heart_failure_dataset.csv
  mall_customers_dataset.csv
  medical_cost_personal_dataset.csv
  optical_recognition_handwritten_digits.zip
  red_wine_quality_dataset.csv
  student_performance_dataset.csv
  titanic_dataset.csv
\end{lstlisting}
\end{minipage}\hfill
\begin{minipage}[t]{0.49\linewidth}
\begin{lstlisting}[style=utpterminal]
all_experiments/
  lab01_workflow.R
  lab02_data_quality.R
  lab03_wrangling.R
  lab04_eda.R
  lab05_visualization.R
  lab06_inference.R
  lab07_regression.R
  lab08_classification.R
  lab09_validation.R
  lab10_unsupervised.R
\end{lstlisting}
\end{minipage}

\section{Route A---recommended: install R natively}

For most students, the simplest and most reliable setup is a native installation of R from the \href{https://cran.r-project.org/}{\textbf{Comprehensive R Archive Network (CRAN)}} followed by RStudio Desktop from \href{https://posit.co/download/rstudio-desktop/}{\textbf{Posit}}.

\subsection{Windows}

\begin{utpsteps}
  \item Open the \href{https://cran.r-project.org/}{CRAN home page}, select \textbf{Download R for Windows}, then choose \textbf{base}.
  \item Download the current stable Windows installer and run it. The default 64-bit installation is suitable for this course.
  \item Keep the default components unless your institution provides different instructions.
  \item After installation, open the Windows Start menu and confirm that an \textbf{R} entry exists.
\end{utpsteps}

From a terminal, installation can be verified with:

\begin{lstlisting}[style=utpterminal]
R --version
\end{lstlisting}

If the terminal cannot find \texttt{R} but R opens normally from the Start menu, the installation itself may still be correct. RStudio can discover standard Windows installations without requiring students to edit the system PATH manually.

\subsection{macOS}

\begin{utpsteps}
  \item Open the \href{https://cran.r-project.org/}{CRAN home page} and select \textbf{Download R for macOS}.
  \item Choose the installer that matches the machine architecture. Recent Apple Silicon Macs use the ARM64 build; older Intel Macs use the Intel build.
  \item Open the downloaded \texttt{.pkg} and follow the installer prompts.
  \item Close and reopen Terminal after installation, then verify R.
\end{utpsteps}

\begin{lstlisting}[style=utpterminal]
R --version
which R
\end{lstlisting}

\subsection{Linux}

Linux installation differs by distribution. Use the distribution-specific instructions under \textbf{Download R for Linux} on CRAN. On Debian/Ubuntu systems, a basic repository installation is commonly performed with:

\begin{lstlisting}[style=utpterminal]
sudo apt update
sudo apt install r-base
\end{lstlisting}

Institutional machines may use a managed repository or module system instead. The objective is the same: the command below should start a current R session.

\begin{lstlisting}[style=utpterminal]
R
\end{lstlisting}

Exit the interactive R console with \cmd{q()} and answer \texttt{n} if asked whether to save the workspace.

\section{Install RStudio Desktop}

RStudio Desktop is available from \href{https://posit.co/download/rstudio-desktop/}{\textbf{Posit's RStudio Desktop download page}}. Posit's current documentation states that RStudio requires an existing R installation, which is why R is installed first.

\begin{description}[leftmargin=1.25in,style=nextline]
  \item[Windows] Download the stable \texttt{.exe}, run the installer, and launch RStudio from the Start menu.
  \item[macOS] Download the stable \texttt{.dmg}, open it, and drag RStudio to \textbf{Applications}.
  \item[Linux] Download the package matching the supported distribution, such as \texttt{.deb} for Ubuntu/Debian or \texttt{.rpm} for Fedora/RHEL-family systems, then install it with the system package manager.
\end{description}

\begin{utpfig}[A useful mental model of the RStudio interface]{The exact pane arrangement can be changed, but the distinction between source code and the console should remain clear. Saved scripts are the reproducible record; console commands are temporary unless copied into a script.}{fig:rstudio-layout}
\begin{tikzpicture}[x=\linewidth,y=1cm,
  pane/.style={draw=UTPBlue!70,fill=UTPPaleBlue,rounded corners=1.5mm,
    line width=0.7pt,text width=0.39\linewidth,minimum height=1.85cm,
    inner xsep=3mm,inner ysep=2mm,align=left,
    font=\footnotesize\color{UTPNavy}}]
\node[pane] (source) at (0,0)
  {\textbf{Source editor}\\[5pt]
   Write and save \texttt{.R} scripts.\\
   Run selected lines or the whole file.};
\node[pane,draw=UTPMuted!70,fill=UTPPaleGray] (env) at (0.49,0)
  {\textbf{Environment / History}\\[5pt]
   Inspect objects in memory.\\
   Review earlier console commands.};
\node[pane,draw=UTPTeal!80,fill=UTPPaleTeal] (console) at (0,-2.15)
  {\textbf{Console}\\[5pt]
   Execute R commands.\\
   Read printed output and errors.};
\node[pane] (tools) at (0.49,-2.15)
  {\textbf{Files / Plots / Packages / Help}\\[5pt]
   Browse files and manage packages.\\
   View figures and documentation.};
\end{tikzpicture}
\end{utpfig}

\subsection{Verify that RStudio is using R}

Open RStudio and run these commands in the Console:

\begin{lstlisting}[style=utpr]
R.version.string
R.home()
getwd()
\end{lstlisting}

The first two lines identify the R installation. The working directory should eventually be the supplied \texttt{DATA SCIENCE/} folder once the course project is opened.

\section{Installing R packages}

The core lab scripts are deliberately written so that they can run with base R on a standard installation. Additional packages are valuable for larger projects, publication graphics, reshaping, and reporting.

Packages from CRAN are installed once and loaded in each new R session when needed:

\begin{lstlisting}[style=utpr]
install.packages(c("ggplot2", "dplyr", "tidyr", "readr", "broom"))

library(ggplot2)
library(dplyr)
\end{lstlisting}

\begin{utpfig}[Install once; load when needed]{Package installation writes software to an R library on the computer. \texttt{library()} attaches an installed package to the current R session.}{fig:r-packages}
\begin{tikzpicture}[node distance=1.1cm]
\node[fcbox,text width=3.8cm] (cran) {CRAN package repository};
\node[fcbox,text width=4.2cm,right=of cran] (lib) {Computer's R library\\\texttt{install.packages()} once};
\node[fcbox,text width=4.2cm,right=of lib,fill=UTPPaleTeal,draw=UTPTeal] (session) {Current R session\\\texttt{library()} when needed};
\draw[fcedge] (cran)--(lib); \draw[fcedge] (lib)--(session);
\end{tikzpicture}
\end{utpfig}

\section{Route B---optional: Anaconda or Miniconda with R}

Some students already work in environments managed by \texttt{conda}. This is a valid optional route, especially when a broader workflow combines multiple languages or command-line tools. However, \textbf{native R + RStudio is the recommended course setup} because it has the fewest moving parts.

Download \href{https://www.anaconda.com/download}{\textbf{Anaconda Distribution}} or follow the \href{https://docs.conda.io/projects/conda/en/stable/user-guide/install/}{\textbf{conda installation guide}} for Miniconda.

\begin{referencebox}
\textbf{Current R channel note.} Anaconda announced that it no longer maintains its own R channel. For current R packages inside conda environments, Anaconda directs users toward \textbf{conda-forge} or a native CRAN installation. For this reason, the commands below request R explicitly from conda-forge rather than relying on Anaconda's legacy R packages.
\end{referencebox}

Create a course-specific environment rather than installing project software into \texttt{base}:

\begin{lstlisting}[style=utpterminal]
conda create -n ds-r conda-forge::r-base conda-forge::r-tidyverse
conda activate ds-r
R --version
R
\end{lstlisting}

Inside the R session, verify the environment's R installation:

\begin{lstlisting}[style=utpr]
R.version.string
R.home()
q()
\end{lstlisting}

Use the same environment on later sessions:

\begin{lstlisting}[style=utpterminal]
conda activate ds-r
R
\end{lstlisting}

\begin{pitfall}
Do not mix several installation strategies casually. A machine can contain a native CRAN R, a conda R, and multiple package libraries at the same time. When troubleshooting, first determine which executable is active with \texttt{where R} on Windows or \texttt{which -a R} on macOS/Linux, then check \cmd{R.home()} inside R.
\end{pitfall}

\subsection{RStudio with a conda-managed R}

RStudio normally discovers standard system installations of R. A conda-managed R lives inside an environment-specific directory, so discovery can be less automatic and varies by operating system. Students who require both RStudio and conda should confirm which R instance RStudio launches by checking \cmd{R.home()} and compare it with the R executable shown by the activated conda environment.

For this course, if RStudio is using the native CRAN R while the terminal uses conda R, choose one route for course work and stay consistent. Reproducibility improves when all labs use the same interpreter and package library.